% Généré par measure_zk_resources.py — le label est définitif, adapter le caption si besoin.
% Fusion des anciennes tables « artefacts » et « size of the data handled by proof
% generation » : l'étiquette de la première colonne porte le rôle de l'artefact dans
% la génération de preuve (input / intermediate / output).
% Setup artefacts (Powers of Tau, proving key) and the witness grow with the circuit,
% whereas the proof and the public inputs stay constant -- the on-chain payload does
% not depend on how many transactions the batch contains.
% Proof size at the largest batch : 805 B.
\begin{table}[t]
\centering
\caption{Pipeline artefacts and size of the data handled by proof generation, as a function of the batch size $N_b$. The tag in the first column gives the role of the artefact in proof generation (input, intermediate witness, output).}
\label{tab:zk_artifacts}
\renewcommand{\arraystretch}{1.2}
\resizebox{\columnwidth}{!}{
\begin{tabular}{|l|l|l|r|}
\hline
\textbf{Artefact} & \textbf{Produced by} & \textbf{Once per} & $N_b = 512$ \\
\hline
R1CS system & Setup (circom) & circuit & 13.6\,MiB \\
Proving key\,\footnotesize(input) & Setup (phase 2) & circuit & 45.5\,MiB \\
Witness\,\footnotesize(intermediate) & Proving & batch (ephemeral) & 2.3\,MiB \\
Groth16 proof\,\footnotesize(output) & Proving & batch (on-chain) & 805\,B \\
Public inputs\,\footnotesize(output) & Proving & batch (on-chain) & 8\,B \\
\hline
\end{tabular}}
\end{table}
